\uuid{oSrW}
\exo7id{530}
\titre{exo7 530}
\auteur{cousquer}
\organisation{exo7}
\datecreate{2003-10-01}
\isIndication{false}
\isCorrection{false}
\chapitre{Suite}
\sousChapitre{Convergence}
\cpge{mpsi-mp2i.nombres-reels-suites.c, pcsi.nombres-reels-suites.b, ptsi.nombres-reels-suites.b, mpsi-mp2i.nombres-reels-suites.d, pcsi.nombres-reels-suites.c, ptsi.nombres-reels-suites.c, mpsi-mp2i.nombres-reels-suites.f, pcsi.nombres-reels-suites.e, ptsi.nombres-reels-suites.e, bcpst1.suites-reelles}
\module{Analyse}
\niveau{L1}
\difficulte{}

\contenu{
\question{
On donne la suite $(u_n)$ définie par~:
$$ u_1= \sqrt2 \quad\mbox{et}\quad u_n = \sqrt{2-u_{n-1}}.$$
En étudiant les suites $(u_{2n})$ et $(u_{2n+1})$,
montrer que la suite $(u_n)$ est convergente.
}
}
